% --- Archon physics formalization source begin ---
% archon:physics
% archon:covers PhyXMiniProblems/problem_phyx_mini_0313.lean
% archon:source-report reports/phyx_mini/problem_phyx_mini_0313.source.json

\chapter{Physics problem phyx\_mini\_0313}
\label{ch:PhyXMiniProblems_problem_phyx_mini_0313}

\paragraph{Problem source.}
\#\# Physical scenario

Figure shows four isotropic point sources of sound that are uniformly spaced on an \$x\$ axis. The sources emit sound at the same wavelength \$\textbackslash{}lambda\$ and same amplitude \$s\_m\$, and they emit in phase. A point \$P\$ is shown on the \$x\$ axis. Assume that as the sound waves travel to \$P\$, the decrease in their amplitude is negligible.

\#\# Question

What multiple of \$s\_m\$ is the amplitude of the net wave at \$P\$ if distance \$d\$ in the figure is \$\textbackslash{}lambda\$?

\#\# Answer choices

- A: A: \$s\_m\$

- B: B: \$2s\_m\$

- C: C: \$3s\_m\$

- D: D: \$4s\_m\$

\#\# Figure caption (auxiliary; use the image as primary evidence)

The image depicts a horizontal line with a series of elements placed along it. Here's a breakdown:

- At the left end, there is a red dot labeled \textbackslash{}( S\_1 \textbackslash{}).
- Moving to the right, there are three more red dots labeled \textbackslash{}( S\_2 \textbackslash{}), \textbackslash{}( S\_3 \textbackslash{}), and \textbackslash{}( S\_4 \textbackslash{}) in sequence.
- Each of these red dots is equally spaced, with the distance between them marked as \textbackslash{}( d \textbackslash{}).
- At the far right, there is a star-shaped icon labeled \textbackslash{}( P \textbackslash{}).
- The sequence from left to right is \textbackslash{}( S\_1 \textbackslash{}), \textbackslash{}( S\_2 \textbackslash{}), \textbackslash{}( S\_3 \textbackslash{}), \textbackslash{}( S\_4 \textbackslash{}), and then \textbackslash{}( P \textbackslash{}).

The image conveys a linear arrangement with consistent spacing between each labeled dot.

\#\# Dataset metadata

Resonance and Harmonics; Physical Model Grounding Reasoning, Spatial Relation Reasoning

\paragraph{Recorded answer/context.}
D: D: \$4s\_m\$

\paragraph{Figure/image path.}
/root/proposal\_for\_physic/science-mango/phyx\_mini\_run/phyx\_data/test\_image/313.png

\paragraph{Formalization target.}
create a compiling Lean file with sorry bodies at `PhyXMiniProblems/problem\_phyx\_mini\_0313.lean`. The Lean declarations must preserve the physical quantities, dimensions or dimensional roles, figure labels, governing-law hypotheses, and final relation expressed by this problem.
Use Mathlib/Physlib names found through LeanExplore where available. If a physics API is missing, introduce faithful local abstractions rather than scalar placeholder aliases.

\begin{theorem}[Physics formalization target]
\label{thm:physics:phyx_mini_0313:target}
\lean{PhyXMiniProblems.ProblemPhyXMini0313.problem_phyx_mini_0313}
\uses{def:physics:phyx-mini-0313:phyxminiproblems-problemphyxmini0313-amplitudereadout, def:physics:phyx-mini-0313:phyxminiproblems-problemphyxmini0313-foursourcesoundsetup, def:physics:phyx-mini-0313:phyxminiproblems-problemphyxmini0313-matchesfoursourcefigure, def:physics:phyx-mini-0313:phyxminiproblems-problemphyxmini0313-hasidenticalcoherentpointsources, def:physics:phyx-mini-0313:phyxminiproblems-problemphyxmini0313-spacingequalswavelength, def:physics:phyx-mini-0313:phyxminiproblems-problemphyxmini0313-hasphysicalparameters, def:physics:phyx-mini-0313:phyxminiproblems-problemphyxmini0313-satisfiescollinearraygeometry, def:physics:phyx-mini-0313:phyxminiproblems-problemphyxmini0313-satisfiesmonochromaticpropagation, def:physics:phyx-mini-0313:phyxminiproblems-problemphyxmini0313-satisfiesnegligibleattenuation, def:physics:phyx-mini-0313:phyxminiproblems-problemphyxmini0313-satisfiescoherentphasorsuperposition, def:physics:phyx-mini-0313:phyxminiproblems-problemphyxmini0313-recordedanswerchoice, def:physics:phyx-mini-0313:phyxminiproblems-problemphyxmini0313-isuniquematchingdisplayedchoice}
The assigned autoformalize agent should translate this physics problem into Lean declarations in the covered file, with theorem and lemma proof bodies written as `by sorry`.
\end{theorem}
\begin{proof}
This is an autoformalization task, not a proof task. Produce statements that can later be proved by the physics prover without weakening the physical model.
\end{proof}

% --- Archon declaration topology begin ---
\section{Declaration topology}
\begin{definition}[Length Quantity]
  \label{def:physics:phyx-mini-0313:phyxminiproblems-problemphyxmini0313-lengthquantity}
  \lean{PhyXMiniProblems.ProblemPhyXMini0313.LengthQuantity}
  A nonnegative physical length such as a wavelength or ray-path length
\end{definition}
\begin{proof}
  This is the declared model definition of Length Quantity.
\end{proof}
\begin{definition}[Signed Length Quantity]
  \label{def:physics:phyx-mini-0313:phyxminiproblems-problemphyxmini0313-signedlengthquantity}
  \lean{PhyXMiniProblems.ProblemPhyXMini0313.SignedLengthQuantity}
  A signed physical coordinate on the horizontal axis in the figure
\end{definition}
\begin{proof}
  This is the declared model definition of Signed Length Quantity.
\end{proof}
\begin{definition}[Acoustic Displacement Amplitude]
  \label{def:physics:phyx-mini-0313:phyxminiproblems-problemphyxmini0313-acousticdisplacementamplitude}
  \lean{PhyXMiniProblems.ProblemPhyXMini0313.AcousticDisplacementAmplitude}
  A nonnegative longitudinal-displacement amplitude of a sound wave
\end{definition}
\begin{proof}
  This is the declared model definition of Acoustic Displacement Amplitude.
\end{proof}
\begin{definition}[length Readout]
  \label{def:physics:phyx-mini-0313:phyxminiproblems-problemphyxmini0313-lengthreadout}
  \lean{PhyXMiniProblems.ProblemPhyXMini0313.lengthReadout}
  \uses{def:physics:phyx-mini-0313:phyxminiproblems-problemphyxmini0313-lengthquantity}
  Read a nonnegative physical length in a selected length unit
\end{definition}
\begin{proof}
  This is the declared model definition of length Readout.
\end{proof}
\begin{definition}[signed Length Readout]
  \label{def:physics:phyx-mini-0313:phyxminiproblems-problemphyxmini0313-signedlengthreadout}
  \lean{PhyXMiniProblems.ProblemPhyXMini0313.signedLengthReadout}
  \uses{def:physics:phyx-mini-0313:phyxminiproblems-problemphyxmini0313-signedlengthquantity}
  Read a signed horizontal coordinate in a selected length unit
\end{definition}
\begin{proof}
  This is the declared model definition of signed Length Readout.
\end{proof}
\begin{definition}[amplitude Readout]
  \label{def:physics:phyx-mini-0313:phyxminiproblems-problemphyxmini0313-amplitudereadout}
  \lean{PhyXMiniProblems.ProblemPhyXMini0313.amplitudeReadout}
  \uses{def:physics:phyx-mini-0313:phyxminiproblems-problemphyxmini0313-acousticdisplacementamplitude}
  Read a sound-displacement amplitude in a selected length unit
\end{definition}
\begin{proof}
  This is the declared model definition of amplitude Readout.
\end{proof}
\begin{definition}[Source Label]
  \label{def:physics:phyx-mini-0313:phyxminiproblems-problemphyxmini0313-sourcelabel}
  \lean{PhyXMiniProblems.ProblemPhyXMini0313.SourceLabel}
  The four source labels printed from left to right in the figure
\end{definition}
\begin{proof}
  This is the declared model definition of Source Label.
\end{proof}
\begin{definition}[Radiation Pattern]
  \label{def:physics:phyx-mini-0313:phyxminiproblems-problemphyxmini0313-radiationpattern}
  \lean{PhyXMiniProblems.ProblemPhyXMini0313.RadiationPattern}
  The angular radiation pattern of a sound source
\end{definition}
\begin{proof}
  This is the declared model definition of Radiation Pattern.
\end{proof}
\begin{definition}[Point Sound Source]
  \label{def:physics:phyx-mini-0313:phyxminiproblems-problemphyxmini0313-pointsoundsource}
  \lean{PhyXMiniProblems.ProblemPhyXMini0313.PointSoundSource}
  \uses{def:physics:phyx-mini-0313:phyxminiproblems-problemphyxmini0313-lengthquantity, def:physics:phyx-mini-0313:phyxminiproblems-problemphyxmini0313-acousticdisplacementamplitude, def:physics:phyx-mini-0313:phyxminiproblems-problemphyxmini0313-radiationpattern}
  A monochromatic point source with physical wavelength and amplitude
\end{definition}
\begin{proof}
  This is the declared model definition of Point Sound Source.
\end{proof}
\begin{definition}[Four Source Sound Setup]
  \label{def:physics:phyx-mini-0313:phyxminiproblems-problemphyxmini0313-foursourcesoundsetup}
  \lean{PhyXMiniProblems.ProblemPhyXMini0313.FourSourceSoundSetup}
  \uses{def:physics:phyx-mini-0313:phyxminiproblems-problemphyxmini0313-lengthquantity, def:physics:phyx-mini-0313:phyxminiproblems-problemphyxmini0313-signedlengthquantity, def:physics:phyx-mini-0313:phyxminiproblems-problemphyxmini0313-acousticdisplacementamplitude, def:physics:phyx-mini-0313:phyxminiproblems-problemphyxmini0313-sourcelabel, def:physics:phyx-mini-0313:phyxminiproblems-problemphyxmini0313-pointsoundsource}
  The independent objects, physical quantities, and observables in the problem. The received amplitudes and phases, and especially netAmplitudeAtP, are left as independent physical observables. They acquire their physical meaning only through the propagation, attenuation, and superposition laws below; no field assigns the requested multiple of s\_m
\end{definition}
\begin{proof}
  This is the declared model definition of Four Source Sound Setup.
\end{proof}
\begin{definition}[Matches Four Source Figure]
  \label{def:physics:phyx-mini-0313:phyxminiproblems-problemphyxmini0313-matchesfoursourcefigure}
  \lean{PhyXMiniProblems.ProblemPhyXMini0313.MatchesFourSourceFigure}
  \uses{def:physics:phyx-mini-0313:phyxminiproblems-problemphyxmini0313-lengthreadout, def:physics:phyx-mini-0313:phyxminiproblems-problemphyxmini0313-signedlengthreadout, def:physics:phyx-mini-0313:phyxminiproblems-problemphyxmini0313-foursourcesoundsetup}
  The one-dimensional layout read from the primary image. The four sources are ordered S₁, S₂, S₃, S₄, point P is to their right, and each adjacent source coordinate differs by the physical spacing d. The image does not label the distance from S₄ to P, so no value is imposed on that distance
\end{definition}
\begin{proof}
  This is the declared model definition of Matches Four Source Figure.
\end{proof}
\begin{definition}[Has Identical Coherent Point Sources]
  \label{def:physics:phyx-mini-0313:phyxminiproblems-problemphyxmini0313-hasidenticalcoherentpointsources}
  \lean{PhyXMiniProblems.ProblemPhyXMini0313.HasIdenticalCoherentPointSources}
  \uses{def:physics:phyx-mini-0313:phyxminiproblems-problemphyxmini0313-sourcelabel, def:physics:phyx-mini-0313:phyxminiproblems-problemphyxmini0313-foursourcesoundsetup}
  The stated source data: all four point sources are isotropic, use the named common wavelength and displacement amplitude, and emit with one initial phase. This does not assert that the waves arrive in phase at P
\end{definition}
\begin{proof}
  This is the declared model definition of Has Identical Coherent Point Sources.
\end{proof}
\begin{definition}[Spacing Equals Wavelength]
  \label{def:physics:phyx-mini-0313:phyxminiproblems-problemphyxmini0313-spacingequalswavelength}
  \lean{PhyXMiniProblems.ProblemPhyXMini0313.SpacingEqualsWavelength}
  \uses{def:physics:phyx-mini-0313:phyxminiproblems-problemphyxmini0313-foursourcesoundsetup}
  The condition in the question, namely that adjacent spacing is d = λ
\end{definition}
\begin{proof}
  This is the declared model definition of Spacing Equals Wavelength.
\end{proof}
\begin{definition}[Has Physical Parameters]
  \label{def:physics:phyx-mini-0313:phyxminiproblems-problemphyxmini0313-hasphysicalparameters}
  \lean{PhyXMiniProblems.ProblemPhyXMini0313.HasPhysicalParameters}
  \uses{def:physics:phyx-mini-0313:phyxminiproblems-problemphyxmini0313-lengthreadout, def:physics:phyx-mini-0313:phyxminiproblems-problemphyxmini0313-amplitudereadout, def:physics:phyx-mini-0313:phyxminiproblems-problemphyxmini0313-sourcelabel, def:physics:phyx-mini-0313:phyxminiproblems-problemphyxmini0313-foursourcesoundsetup}
  Positivity and nondegeneracy of the physical lengths and amplitudes
\end{definition}
\begin{proof}
  This is the declared model definition of Has Physical Parameters.
\end{proof}
\begin{definition}[Satisfies Collinear Ray Geometry]
  \label{def:physics:phyx-mini-0313:phyxminiproblems-problemphyxmini0313-satisfiescollinearraygeometry}
  \lean{PhyXMiniProblems.ProblemPhyXMini0313.SatisfiesCollinearRayGeometry}
  \uses{def:physics:phyx-mini-0313:phyxminiproblems-problemphyxmini0313-lengthreadout, def:physics:phyx-mini-0313:phyxminiproblems-problemphyxmini0313-signedlengthreadout, def:physics:phyx-mini-0313:phyxminiproblems-problemphyxmini0313-sourcelabel, def:physics:phyx-mini-0313:phyxminiproblems-problemphyxmini0313-foursourcesoundsetup}
  Because all objects lie on the same axis and P is to the right, each ray length is the coordinate of P minus the corresponding source coordinate
\end{definition}
\begin{proof}
  This is the declared model definition of Satisfies Collinear Ray Geometry.
\end{proof}
\begin{definition}[Satisfies Monochromatic Propagation]
  \label{def:physics:phyx-mini-0313:phyxminiproblems-problemphyxmini0313-satisfiesmonochromaticpropagation}
  \lean{PhyXMiniProblems.ProblemPhyXMini0313.SatisfiesMonochromaticPropagation}
  \uses{def:physics:phyx-mini-0313:phyxminiproblems-problemphyxmini0313-lengthreadout, def:physics:phyx-mini-0313:phyxminiproblems-problemphyxmini0313-sourcelabel, def:physics:phyx-mini-0313:phyxminiproblems-problemphyxmini0313-foursourcesoundsetup}
  Monochromatic propagation subtracts 2π r / λ from the source phase. The phase is an element of Real.Angle, so paths differing by one wavelength produce the same arrival phase
\end{definition}
\begin{proof}
  This is the declared model definition of Satisfies Monochromatic Propagation.
\end{proof}
\begin{definition}[Satisfies Negligible Attenuation]
  \label{def:physics:phyx-mini-0313:phyxminiproblems-problemphyxmini0313-satisfiesnegligibleattenuation}
  \lean{PhyXMiniProblems.ProblemPhyXMini0313.SatisfiesNegligibleAttenuation}
  \uses{def:physics:phyx-mini-0313:phyxminiproblems-problemphyxmini0313-amplitudereadout, def:physics:phyx-mini-0313:phyxminiproblems-problemphyxmini0313-sourcelabel, def:physics:phyx-mini-0313:phyxminiproblems-problemphyxmini0313-foursourcesoundsetup}
  The stated negligible-decrease approximation: propagation does not change the displacement amplitude of any one source before it reaches P
\end{definition}
\begin{proof}
  This is the declared model definition of Satisfies Negligible Attenuation.
\end{proof}
\begin{definition}[Satisfies Coherent Phasor Superposition]
  \label{def:physics:phyx-mini-0313:phyxminiproblems-problemphyxmini0313-satisfiescoherentphasorsuperposition}
  \lean{PhyXMiniProblems.ProblemPhyXMini0313.SatisfiesCoherentPhasorSuperposition}
  \uses{def:physics:phyx-mini-0313:phyxminiproblems-problemphyxmini0313-amplitudereadout, def:physics:phyx-mini-0313:phyxminiproblems-problemphyxmini0313-sourcelabel, def:physics:phyx-mini-0313:phyxminiproblems-problemphyxmini0313-foursourcesoundsetup}
  General coherent superposition at P. The net amplitude is the magnitude of the sum of the four received phasors. This law is uniform in the received phases and amplitudes and contains neither the value 4 s\_m nor an answer choice
\end{definition}
\begin{proof}
  This is the declared model definition of Satisfies Coherent Phasor Superposition.
\end{proof}
\begin{lemma}[ray Path Differences Are Whole Wavelengths]
  \label{lem:physics:phyx-mini-0313:phyxminiproblems-problemphyxmini0313-raypathdifferencesarewholewavelengths}
  \lean{PhyXMiniProblems.ProblemPhyXMini0313.rayPathDifferencesAreWholeWavelengths}
  \uses{def:physics:phyx-mini-0313:phyxminiproblems-problemphyxmini0313-lengthreadout, def:physics:phyx-mini-0313:phyxminiproblems-problemphyxmini0313-sourcelabel, def:physics:phyx-mini-0313:phyxminiproblems-problemphyxmini0313-foursourcesoundsetup, def:physics:phyx-mini-0313:phyxminiproblems-problemphyxmini0313-matchesfoursourcefigure, def:physics:phyx-mini-0313:phyxminiproblems-problemphyxmini0313-spacingequalswavelength, def:physics:phyx-mini-0313:phyxminiproblems-problemphyxmini0313-satisfiescollinearraygeometry}
  The depicted equal spacing together with d = λ makes every source-to-P path differ from the S₄ path by an integral number of wavelengths
\end{lemma}
\begin{proof}
  The conclusion follows from the governing relations and figure readouts named in the statement of ray Path Differences Are Whole Wavelengths.
\end{proof}
\begin{lemma}[all Received Phases Agree]
  \label{lem:physics:phyx-mini-0313:phyxminiproblems-problemphyxmini0313-allreceivedphasesagree}
  \lean{PhyXMiniProblems.ProblemPhyXMini0313.allReceivedPhasesAgree}
  \uses{def:physics:phyx-mini-0313:phyxminiproblems-problemphyxmini0313-sourcelabel, def:physics:phyx-mini-0313:phyxminiproblems-problemphyxmini0313-foursourcesoundsetup, def:physics:phyx-mini-0313:phyxminiproblems-problemphyxmini0313-matchesfoursourcefigure, def:physics:phyx-mini-0313:phyxminiproblems-problemphyxmini0313-hasidenticalcoherentpointsources, def:physics:phyx-mini-0313:phyxminiproblems-problemphyxmini0313-spacingequalswavelength, def:physics:phyx-mini-0313:phyxminiproblems-problemphyxmini0313-hasphysicalparameters, def:physics:phyx-mini-0313:phyxminiproblems-problemphyxmini0313-satisfiescollinearraygeometry, def:physics:phyx-mini-0313:phyxminiproblems-problemphyxmini0313-satisfiesmonochromaticpropagation}
  Equal source phases and integral-wavelength path differences imply that all four received waves have the same phase at P
\end{lemma}
\begin{proof}
  The conclusion follows from the governing relations and figure readouts named in the statement of all Received Phases Agree.
\end{proof}
\begin{lemma}[coherent Phasor Components]
  \label{lem:physics:phyx-mini-0313:phyxminiproblems-problemphyxmini0313-coherentphasorcomponents}
  \lean{PhyXMiniProblems.ProblemPhyXMini0313.coherentPhasorComponents}
  \uses{def:physics:phyx-mini-0313:phyxminiproblems-problemphyxmini0313-amplitudereadout, def:physics:phyx-mini-0313:phyxminiproblems-problemphyxmini0313-sourcelabel, def:physics:phyx-mini-0313:phyxminiproblems-problemphyxmini0313-foursourcesoundsetup, def:physics:phyx-mini-0313:phyxminiproblems-problemphyxmini0313-hasidenticalcoherentpointsources, def:physics:phyx-mini-0313:phyxminiproblems-problemphyxmini0313-satisfiesnegligibleattenuation}
  With negligible attenuation and equal arrival phases, the Cartesian phasor components are four times the corresponding component of one source
\end{lemma}
\begin{proof}
  The conclusion follows from the governing relations and figure readouts named in the statement of coherent Phasor Components.
\end{proof}
\begin{definition}[Answer Choice]
  \label{def:physics:phyx-mini-0313:phyxminiproblems-problemphyxmini0313-answerchoice}
  \lean{PhyXMiniProblems.ProblemPhyXMini0313.AnswerChoice}
  The four answer labels printed with the problem
\end{definition}
\begin{proof}
  This is the declared model definition of Answer Choice.
\end{proof}
\begin{definition}[amplitude Multiplier]
  \label{def:physics:phyx-mini-0313:phyxminiproblems-problemphyxmini0313-answerchoice-amplitudemultiplier}
  \lean{PhyXMiniProblems.ProblemPhyXMini0313.AnswerChoice.amplitudeMultiplier}
  \uses{def:physics:phyx-mini-0313:phyxminiproblems-problemphyxmini0313-answerchoice}
  The multiple of s\_m printed beside each answer label
\end{definition}
\begin{proof}
  This is the declared model definition of amplitude Multiplier.
\end{proof}
\begin{definition}[recorded Answer Choice]
  \label{def:physics:phyx-mini-0313:phyxminiproblems-problemphyxmini0313-recordedanswerchoice}
  \lean{PhyXMiniProblems.ProblemPhyXMini0313.recordedAnswerChoice}
  \uses{def:physics:phyx-mini-0313:phyxminiproblems-problemphyxmini0313-answerchoice}
  The answer label recorded in the supplied dataset metadata
\end{definition}
\begin{proof}
  This is the declared model definition of recorded Answer Choice.
\end{proof}
\begin{definition}[Matches Displayed Amplitude Choice]
  \label{def:physics:phyx-mini-0313:phyxminiproblems-problemphyxmini0313-matchesdisplayedamplitudechoice}
  \lean{PhyXMiniProblems.ProblemPhyXMini0313.MatchesDisplayedAmplitudeChoice}
  \uses{def:physics:phyx-mini-0313:phyxminiproblems-problemphyxmini0313-amplitudereadout, def:physics:phyx-mini-0313:phyxminiproblems-problemphyxmini0313-foursourcesoundsetup, def:physics:phyx-mini-0313:phyxminiproblems-problemphyxmini0313-answerchoice}
  A displayed choice matches when its multiplier gives the physical net amplitude in every length unit. This generic definition does not select D
\end{definition}
\begin{proof}
  This is the declared model definition of Matches Displayed Amplitude Choice.
\end{proof}
\begin{definition}[Is Unique Matching Displayed Choice]
  \label{def:physics:phyx-mini-0313:phyxminiproblems-problemphyxmini0313-isuniquematchingdisplayedchoice}
  \lean{PhyXMiniProblems.ProblemPhyXMini0313.IsUniqueMatchingDisplayedChoice}
  \uses{def:physics:phyx-mini-0313:phyxminiproblems-problemphyxmini0313-foursourcesoundsetup, def:physics:phyx-mini-0313:phyxminiproblems-problemphyxmini0313-answerchoice, def:physics:phyx-mini-0313:phyxminiproblems-problemphyxmini0313-matchesdisplayedamplitudechoice}
  A displayed multiplier is the unique one matching the physical result
\end{definition}
\begin{proof}
  This is the declared model definition of Is Unique Matching Displayed Choice.
\end{proof}
% --- Archon declaration topology end ---
% --- Archon physics formalization source end ---
